\section*{Capítulo \ref{ch:b2:ism-star-formation} --- O meio interestelar e a formação das estrelas}

\begin{solution}{exo:b2:ism-star-formation:1}
$n = \num{2.7e-24}/\num{1.67e-24} = 1.6$ átomos de hidrogênio por centímetro cúbico. No ar,
$n = P/k_{\mathrm B}T = 101325/(\num{1.381e-23}\times288) = \qty{2.55e25}{m^{-3}}$,
isto é, \num{2.5e19} moléculas por centímetro cúbico: cerca de $1.6\times10^{19}$ vezes
mais.
\end{solution}

\begin{solution}{exo:b2:ism-star-formation:2}
$M_J \propto n^{-1/2}$: dividir $n$ por 100 multiplica $M_J$ por 10, de modo que
$M_J = 54.6$ massas solares com $10^2$ partículas por centímetro cúbico. A nuvem contém
$\num{24000}/54.6 \approx 440$ \omterm{def:b2:ism-star-formation:jeans}{massas de Jeans}: ela não pode ser estável como um todo e precisa
se fragmentar.
\end{solution}

\begin{solution}{exo:b2:ism-star-formation:3}
$c_s = \sqrt{\num{1.381e-23}\times10/(2.33\times\num{1.661e-27})} =
\qty{189}{m.s^{-1}}$ a \qty{10}{K}; $\sqrt{10}$ vezes mais, \qty{597}{m.s^{-1}},
a \qty{100}{K}. Travessia de \qty{0.1}{pc}: $\num{3.086e15}/189 = \qty{1.63e13}{s}$,
isto é, \qty{0.52}{Myr}.
\end{solution}

\begin{solution}{exo:b2:ism-star-formation:4}
$\rho = 2.33\times\num{1.661e-27}\times n$ com $n$ em \unit{m^{-3}}. Para $10^5$
partículas por centímetro cúbico, $\rho = \qty{3.87e-16}{kg.m^{-3}}$ e
$t_{\mathrm{ff}} = \sqrt{3\pi/(32G\rho)} = \qty{3.38e12}{s} = \qty{0.107}{Myr}$; para
$10^2$, $\rho$ é \num{1000} vezes menor e $t_{\mathrm{ff}}$, $\sqrt{1000} = 31.6$
vezes mais longo, \qty{3.38}{Myr}. O \omterm{def:b2:ism-star-formation:molecular-cloud}{núcleo denso} colapsa 32 vezes mais depressa.
\end{solution}

\begin{solution}{exo:b2:ism-star-formation:5}
Com $k_{\mathrm B}T/(\mu m_u) = c_s^2$, $(5k_{\mathrm B}T/G\mu m_u)^{3/2} =
5^{3/2}c_s^3G^{-3/2}$, o que dá a fórmula. Como $M_J \propto \rho^{-1/2}$,
reduzir $M_J$ à metade exige uma densidade quatro vezes maior. Como $M_J \propto T^{3/2}$,
aquecer de \qty{10}{K} para \qty{20}{K} a multiplica por $2^{3/2} = 2.83$.
\end{solution}

\begin{solution}{exo:b2:ism-star-formation:6}
$t_{\mathrm{KH}} = GM^2/(RL) = \num{6.674e-11}\times(\num{1.989e30})^2/(2\times
\num{6.96e8}\times2\times\num{3.83e26}) = \qty{2.48e14}{s}$, isto é,
\qty{7.8}{Myr}. É mais longo que a duração de vida da maioria dos discos (metade desaparece depois de
cerca de \qty{3}{Myr}): a estrela ainda está se contraindo quando seu disco desaparece.
\end{solution}

\begin{solution}{exo:b2:ism-star-formation:7}
A \qty{10}{K}, $c_s^3/G = 189^3/\num{6.674e-11} = \qty{1.01e17}{kg.s^{-1}}$, isto é,
$\num{1.01e17}\times\num{3.156e7}/\num{1.989e30} = \num{1.6e-6}$ massas solares por
ano, e uma massa solar leva \qty{0.62}{Myr}. A \qty{20}{K}, a taxa é
$2^{3/2} = 2.83$ vezes maior, \num{4.5e-6} massas solares por ano, e uma massa
solar leva \qty{0.22}{Myr}.
\end{solution}

\begin{solution}{exo:b2:ism-star-formation:8}
$r = \sqrt{L/(16\pi\sigma)}/T^2$ com $\sqrt{\num{3.83e26}/(16\pi\times\num{5.67e-8})}
= \qty{1.159e16}{m.K^2}$: a \qty{1300}{K}, $r = \qty{6.86e9}{m} = \qty{0.046}{au}$,
bem dentro da órbita de Mercúrio. A \qty{5.2}{au}, $T = 278/\sqrt{5.2} = \qty{122}{K}$.
\end{solution}

\begin{solution}{exo:b2:ism-star-formation:9}
$v = \sqrt{GM/r} = \sqrt{\num{6.674e-11}\times1.3\times\num{1.989e30}/\num{1.496e13}}
= \qty{3.40}{km.s^{-1}}$, e $P = 2\pi r/v = \qty{2.77e10}{s} = 877$ anos
(ou, de modo equivalente, $\sqrt{100^3/1.3}$ anos). Em um milhão de anos, o gás deu
$10^6/877 \approx \num{1140}$ voltas: o disco é velho comparado com sua própria rotação,
e é por isso que seus anéis são círculos regulares.
\end{solution}

\begin{solution}{exo:b2:ism-star-formation:10}
$M = \int_0^R 4\pi r^2\rho_0R/r\,\mathrm dr = 2\pi\rho_0R^3$, e, dentro de $r$,
$m(r) = 2\pi\rho_0Rr^2$. Com $\mathrm dm = 4\pi\rho_0Rr\,\mathrm dr$,
\[
  U = -\int_0^R \frac{G\,m(r)}{r}\,\mathrm dm
    = -8\pi^2G\rho_0^2R^2\int_0^R r^2\,\mathrm dr = -\frac{8\pi^2}{3}G\rho_0^2R^5
    = -\frac23\,\frac{GM^2}{R},
\]
usando $\rho_0 = M/(2\pi R^3)$. Ela vale $\tfrac{2/3}{3/5} = 1.11$ vezes o valor homogêneo:
a nuvem concentrada é mais fortemente ligada e precisa de \qty{11}{\percent} mais
energia térmica para ser sustentada; ela colapsa mais facilmente.
\end{solution}

\begin{solution}{exo:b2:ism-star-formation:11}
Período $2\pi/\Omega = \qty{6.28e14}{s}$, cerca de \qty{20}{Myr}. Com
$R = \qty{3.086e15}{m}$, $j = \Omega R^2 = \qty{9.52e16}{m^2.s^{-1}}$, e
$r_c = j^2/(GM) = \num{9.07e33}/\num{1.327e20} = \qty{6.83e13}{m} = \qty{457}{au}$.
Em rotação de corpo rígido, $j = \Omega r^2$ cresce com o quadrado da distância ao eixo:
o gás do interior do núcleo tem menos momento angular e cai mais perto da estrela.
\end{solution}

\begin{solution}{exo:b2:ism-star-formation:12}
$\bar t = 2.46$, $\bar f = 0.588$; $\sum(t - \bar t)^2 = 12.612$,
$\sum(t - \bar t)(f - \bar f) = -2.018$, portanto $b = 2.018/12.612 = 0.160$ por milhão
de anos e $a = 0.588 + 0.160\times2.46 = 0.982$. A reta atinge $0.5$ em
$(0.982 - 0.5)/0.160 = \qty{3.0}{Myr}$ e zero em $0.982/0.160 = \qty{6.1}{Myr}$, como na
figura. Uma fração não pode cair abaixo de zero: aglomerados mais velhos, como NGC 1960, com
\qty{30}{Myr}, conservam uma pequena fração próxima de zero que uma reta não consegue descrever.
\end{solution}

\begin{solution}{pb:b2:ism-star-formation:1}
\textbf{1.} $R = 0.08\times\num{3.086e16} = \qty{2.469e15}{m}$.

\textbf{2.} $N = \qty{1.5e26}{m^{-2}}$, portanto $2.8\times\num{1.661e-27}\times\num{1.5e26} =
\qty{0.698}{kg.m^{-2}}$.

\textbf{3.} $M = 0.698\times\pi R^2 = 0.698\times\num{1.915e31} = \qty{1.34e31}{kg}$,
isto é, $6.7$ massas solares.

\textbf{4.} $V = \tfrac43\pi R^3 = \qty{6.30e46}{m^3}$ e
$\rho = M/V = \qty{2.12e-16}{kg.m^{-3}}$.

\textbf{5.} $n = \rho/(2.33m_u) = \qty{5.48e10}{m^{-3}}$, isto é, \num{5.5e4} partículas
por centímetro cúbico.

\textbf{6.} Sim: sua densidade fica entre $10^4$ e $10^5$, seu raio entre 0.03 e
\qty{0.2}{pc}; sua massa, 6.7 massas solares, está no topo da faixa dos núcleos.

\textbf{7.} $M_J = 5.46\times(10^4/\num{5.48e4})^{1/2} = 2.33$ massas solares.

\textbf{8.} $6.7/2.33 = 2.9$ \omterm{def:b2:ism-star-formation:jeans}{massas de Jeans}: o núcleo é instável e colapsa.

\textbf{9.} $R_J = 0.1316\times(10^4/\num{5.48e4})^{1/2} = \qty{0.056}{pc}$, menor que
o raio, \qty{0.08}{pc}: a mesma conclusão.

\textbf{10.} $2^{3/2}\times2.33 = 6.6$ massas solares: aproximadamente a massa do núcleo. A
\qty{20}{K}, o núcleo seria marginal; é o frio que o faz colapsar.

\textbf{11.} $t_{\mathrm{ff}} = \sqrt{3\pi/(32\times\num{6.674e-11}\times\num{2.12e-16})}
= \qty{4.56e12}{s}$, cerca de \num{145000} anos.

\textbf{12.} $R/c_s = \num{2.469e15}/189 = \qty{1.31e13}{s}$, \num{414000} anos: três
vezes mais que a queda livre. Uma onda de pressão não consegue atravessar o núcleo antes que ele
tenha colapsado; a pressão não consegue reagir depressa o bastante para deter o colapso.

\textbf{13.} $c_s = \qty{189}{m.s^{-1}}$ (\cref{exo:b2:ism-star-formation:3}).

\textbf{14.} $c_s^3/G = \qty{1.01e17}{kg.s^{-1}}$, \num{1.6e-6} massas solares por ano.

\textbf{15.} $1/\num{1.6e-6} = \num{624000}$ anos.

\textbf{16.} Um terço de 6.7 é $2.2$: duas estrelas de uma massa solar.

\textbf{17.} $t_{\mathrm{KH}} = GM_\odot^2/(3R_\odot\times4L_\odot) = \qty{8.25e13}{s}$,
\qty{2.6}{Myr}.

\textbf{18.} Mais longo, quatro vezes o tempo de acreção: quando a acreção cessa, a estrela
ainda está se contraindo, brilhando sem fusão; é uma \omterm{def:b2:ism-star-formation:t-tauri}{estrela T Tauri}, ainda fora da
sequência principal.

\textbf{19.} $\Omega = \qty{0.4e3}{m.s^{-1}}/\qty{3.086e16}{m} = \qty{1.30e-14}{rad.s^{-1}}$.

\textbf{20.} $j = 1.30\times10^{-14}\times(\num{2.469e15})^2 = \qty{7.90e16}{m^2.s^{-1}}$.

\textbf{21.} $r_c = j^2/(GM_\odot) = \num{6.24e33}/\num{1.327e20} = \qty{4.70e13}{m}$, isto
é, \qty{314}{au}: maior que o disco de HL Tauri, como se espera de um gás vindo da borda do
núcleo.

\textbf{22.} $T = (L/16\pi\sigma r^2)^{1/4} = \qty{278}{K}$ a \qty{1}{au}, e
$278/\sqrt{30} = \qty{51}{K}$ a \qty{30}{au}.

\textbf{23.} $r = \num{1.159e16}/180^2 = \qty{3.58e11}{m}$, isto é, \qty{2.4}{au}, como no texto.

\textbf{24.} A estrela de uma massa solar leva
\textbf{\qty[text-series-to-math=true]{624}{} mil anos} para se formar, cerca de quatro
vezes o \omterm{def:b2:ism-star-formation:free-fall}{tempo de queda livre} de seu núcleo (\num{145000} anos): o gás exterior chega à
\omterm{def:b2:ism-star-formation:protostar}{protoestrela} só aos poucos, de dentro para fora.
\end{solution}
